\section{Implementation and Validation Notes}
\label{app:implementation}

The kernels specialize the production tensor shapes and define buffer
lifetimes across eager execution, recomputation, and CUDA Graph replay. QKV and
FFN producers emit both quantized orientations. Residual--RMSNorm preserves the
BF16 residual and exact row statistics required by backward. Each path is
checked against a higher-precision reference for outputs, input gradients,
weight gradients, relative RMS error, and zero fractions before model timing.

Row- and column-oriented scales are distinct numerical state, not
interchangeable layout metadata. Tensors returned to autograd cannot alias a
workspace whose lifetime ends with the producing stream. Delayed FP4 prefetch
has one owner and one release point. Distributed measurements use per-rank
timing and dispatch evidence because the slowest rank gates the step.

\subsection{Weight orientation and the fixed sign pattern}

The global-v5 path creates one 2D16 numerical weight
encoding and exposes it in both consumer layouts. The native 2D32 \mxfp{} path
follows the same orientation-consistency rule.

The signed-Hadamard paths use the fixed literal \texttt{0x00002817}. Under the
custom kernel's convention, a set bit denotes $+1$ and a clear bit denotes
$-1$. Bits 0--15 encode the same 16-entry sign vector that the pinned TE
implementation stores explicitly. The H32 implementation operates as two
16-value thread halves, and each half reuses this low-16-bit motif before the
cross-half shuffle. Thus H32 names the 32-value transform width, not a 32-bit
random mask: the reported H32 transform repeats one signed-H16 diagonal and
does not consume the literal's upper zero bits as positions 16--31. No signs
are drawn during training. The construction remains
orthogonal and preserves the unquantized Wgrad dot product because the
identical transform is paired on $X$ and $dY$. The experiments do not compare
this fixed vector with newly sampled signs and do not establish that the
particular diagonal is optimal.

Compilation and first-step kernel generation are excluded from steady-state
probe windows. Each measurement records device count, local batch, gradient
accumulation, seed, loss implementation, source revision, and metric window.
Each stochastic producer reserves one logical random-number interval, and the
ranked seed/subsequence state is checkpointed across resumes.

Table~\ref{tab:b300-snapshots} gives secondary production measurements. All
three use a BF16 output projection and compiled regular CE. Their batch
geometries and measurement windows differ, so they are not a one-factor
comparison.

\begin{table}[H]
  \centering
  \caption{Compiled-CE B300 production snapshots.
  BF16 MFU uses the B300's 2,250 TFLOP/s/GPU dense BF16 peak
  \cite{nvidia2026hgx}.}
  \label{tab:b300-snapshots}
  \begin{tabularx}{\textwidth}{Ylrrl}
    \toprule
    Route & LB/GA & Tokens/s/GPU & BF16 MFU & HBM/GPU \\
    \midrule
    \mxfp{} + row-SR & 4/4 & 39,306.59 & 101.17\% & 156.91--157.15 GiB \\
    \localcta{} & 8/2 & 36,728.43 & 94.54\% & 211.05--211.16 GiB \\
    27/5 depth hybrid & 8/2 & 37,394.00 & 96.25\% & 228.28--233.42 GiB \\
    \bottomrule
  \end{tabularx}
\end{table}
